\documentclass[11pt]{article}

\usepackage[T1]{fontenc}
\usepackage{lmodern}
\usepackage{amsmath,amssymb,amsthm,mathtools}
\usepackage[a4paper,margin=29mm]{geometry}
\usepackage{booktabs}
\usepackage{array}
\usepackage{enumitem}
\usepackage{xcolor}
\usepackage{microtype}
\usepackage[hidelinks]{hyperref}
\usepackage{url}

\hypersetup{
  pdftitle={Practical Public-Key-Only Forgery Against the Submitted VDOO-128 Implementation},
  pdfauthor={Martin Feussner},
  pdfsubject={Cryptanalysis of the submitted VDOO-128 reference implementation},
  pdfkeywords={VDOO, multivariate cryptography, digital signatures, structural cryptanalysis, forgery}
}

\newtheorem{theorem}{Theorem}
\newtheorem{lemma}[theorem]{Lemma}
\newtheorem{proposition}[theorem]{Proposition}
\newtheorem{corollary}[theorem]{Corollary}
\newtheorem{definition}[theorem]{Definition}
\newtheorem{remark}[theorem]{Remark}

\newcommand{\F}{\mathbb{F}}
\newcommand{\Proj}{\mathbb{P}}
\newcommand{\Span}{\operatorname{span}}
\newcommand{\rank}{\operatorname{rank}}
\newcommand{\rad}{\operatorname{rad}}
\newcommand{\essrank}{\operatorname{erank}}
\newcommand{\im}{\operatorname{im}}
\newcommand{\kerof}{\operatorname{ker}}
\newcommand{\defeq}{\mathrel{:=}}
\newcommand{\code}[1]{\texttt{#1}}

\title{\bfseries Practical Public-Key-Only Forgery Against the\\
Submitted VDOO-128 Implementation}
\author{Martin Feussner\\
\small Selmer Center, University of Bergen\\
\small \texttt{martin.feussner@uib.no}}
\date{25 September 2026}

\begin{document}
\maketitle

\begin{abstract}
We give a key-only existential forgery attack against the VDOO-128 reference
implementation submitted to the Next-generation Commercial Cryptographic
Algorithms Program.  The attack is mathematical: it uses only the public
quadratic map, makes no signing queries, and does not use the independently
reported random-number-generator failure.  The central-map generator retains,
in each oil equation, products with only the oil variable having the same
index as that equation.  The submitted map is therefore a hidden
one-variable-per-equation triangular system rather than the dense UOV layers
defined in the specification.

We recover the two hidden oil spaces through rank-one polar derivatives.  A
small projective multipencil search finds all 40 lines of the second oil layer.
After quotienting by their input and output spans, a five-evaluation
meet-in-the-middle multipencil finds the 14 lines of the first oil layer.  A
low-essential-rank filtration then recovers an equivalent flag for the 16
diagonal equations.  This equivalent structure suffices to invert the public
map: first invert the 16-layer quotient, then lift through two linear systems.

For the claimed 128-bit parameter set over $\F_{16}$, an unoptimized
single-threaded Python/NumPy implementation completes a fresh public-key attack
and produces an 85-byte signature accepted by the submitted verifier in under
twelve minutes on a 2015 desktop-class Intel Core i5-6500.  A second experiment
uses an independently seeded key to separate the structural attack from the
submission's deterministic-RNG defect.  The result applies to the submitted
coefficient mask.  It does not establish an attack on the dense oil layers in
the normative VDOO specification.
\end{abstract}

\section{Introduction}

VDOO is a multivariate signature design whose central map combines a diagonal
layer with two Unbalanced Oil-and-Vinegar (UOV) layers
\cite{vdoo-original,vdoo-spec}.  For the NGCC VDOO-128 parameter set, the
public key is a homogeneous quadratic map
\[
       P:\F_{16}^{137}\longrightarrow\F_{16}^{70},
\]
hidden by invertible linear input and output transformations.  Its parameter
tuple is
\[
       (q,v,d,o_1,o_2)=(16,67,16,14,40).
\]
The submission claims 128-bit classical security and reports an 85-byte
signature.

The reference implementation does not generate the oil layers described by
the specification.  In a dense UOV layer, every output polynomial contains
vinegar--oil products with every oil variable.  The submitted coefficient
filter instead gives output $j$ products with only oil variable $j$.  Thus the
last 40 outputs have the form
\[
       h_j(u,z)=a_j(u)+z_j\ell_j(u),\qquad 1\leq j\leq 40,
\]
and, after removing this layer, the preceding 14 outputs have the same form.
The oil linear systems have silently become diagonal systems.  More
significantly for cryptanalysis, every coordinate oil direction has a
rank-one public derivative.

This yields a practical attack with the following stages.
\begin{enumerate}[leftmargin=*,itemsep=2pt]
  \item Recover the 40 hidden second-oil coordinate lines from rank-one
        derivatives and form their input and output spans $O_2,Q_2$.
  \item Pass to the induced map
        $\F_{16}^{137}/O_2\rightarrow\F_{16}^{70}/Q_2$ and recover the 14
        first-oil lines, giving $O_1/O_2$ and $Q_1/Q_2$.
  \item Quotient once more.  The remaining map
        $\F_{16}^{83}\rightarrow\F_{16}^{16}$ is a hidden stepwise
        triangular map.  Recover its nested output and input flags from the
        essential ranks and polar radicals of public quadratic forms.
  \item Invert the 16-layer quotient and lift the preimage through $O_1/O_2$
        and $O_2$ by linear algebra.
\end{enumerate}

The attack belongs to the broad family of attacks on stepwise triangular
systems that use chains of kernels \cite{wolf-sts}.  Its first two stages are
specialized to the submitted VDOO mask: projective multipencils turn the
rank-one derivative signal into two small exhaustive searches.  The attack
recovers only the equivalent structure needed for inversion; it never needs
the submitted secret key.

The public NGCC report already records unrelated implementation failures,
including an uninitialized deterministic generator, prehash limits, and
secret-dependent branches \cite{ngcc-report}.  Our experiments explicitly
initialize the generator for an independent key and discard the secret key
before attacking.  Consequently, fixing the RNG does not remove this attack.

\paragraph{Scope.}
We analyze the VDOO-128 reference source at NGCC harness commit
\code{32f381e}.  The normative formulas in
the VDOO specification contain the full vinegar--oil sums.  Our result is a
complete forgery of the submitted implementation distribution, not a claim
against that dense specified distribution.

\section{Preliminaries}

Let $K=\F_{16}$.  All maps in this paper are over $K$.  For a homogeneous
quadratic map $P:K^n\to K^m$, define its polar map and the derivative at $x$ by
\begin{align*}
 B_P(x,y) &\defeq P(x+y)-P(x)-P(y),\\
 J_x(y)   &\defeq B_P(x,y).
\end{align*}
In characteristic two, subtraction equals addition.  The map $B_P$ is
bilinear and alternating in each scalar quadratic component, while $J_x$ is a
linear map $K^n\to K^m$.

Suppose $P=S\circ F\circ T$ for invertible linear maps $T$ and $S$.  Then
\[
       J_x^P=S\,J_{Tx}^F\,T,
\]
so $\rank J_x$ is invariant under the secret hiding transformations.  The
image line of a rank-one derivative is transformed by $S$, and its input line
is transformed by $T^{-1}$.  These are therefore public invariants.

\subsection{Quadratic rank in characteristic two}

A scalar quadratic form $f$ is represented by its alternating polar matrix
$A_f$ together with its square-term coefficients.  We use the \emph{essential
rank} $\essrank(f)$, the least number of variables needed to express $f$ after
an invertible linear change of variables.  It is computable without solving a
polynomial system:
\[
 \essrank(f)=
 \begin{cases}
   \rank A_f,   & f|_{\rad(A_f)}=0,\\
   \rank A_f+1, & f|_{\rad(A_f)}\neq0.
 \end{cases}
\]
Here $\rad(A_f)=\ker A_f$.  Evaluation of $f$ on a basis of the radical
decides which case holds.  This correction is necessary over characteristic
two; polar rank alone loses square terms.

\subsection{Compatible quotients}

The following elementary observation is used twice.

\begin{lemma}[Quadratic quotient]
Let $Z\subseteq K^n$ and $R\subseteq K^m$ satisfy
\[
        P(Z)=0,\qquad B_P(K^n,Z)\subseteq R.
\]
Then $\overline P(x+Z)=P(x)+R$ is a well-defined quadratic map
$K^n/Z\to K^m/R$.
\end{lemma}

\begin{proof}
For $z\in Z$,
$P(x+z)=P(x)+P(z)+B_P(x,z)\equiv P(x)\pmod R$.
\end{proof}

Once bases for $Z$ and $R$ are public, the quotient is computed with ordinary
Gaussian elimination: select a complement of $Z$, and left-compose the
outputs with a basis of the annihilator of $R$.

\section{The submitted coefficient-mask defect}

Write the central variables as
\[
     (u,d,y,z)\in K^{67}\times K^{16}\times K^{14}\times K^{40}.
\]
The specification defines each first-oil output with all products
$x_i y_j$, for $1\leq i\leq83$ and $1\leq j\leq14$, and each second-oil
output with all products $x_i z_j$, for $1\leq i\leq97$ and
$1\leq j\leq40$ \cite{vdoo-spec}.

In the submitted source, the admissibility predicate for first-oil output
$r$ sets
\[
       \text{oil\_var\_idx}=83+r
\]
and accepts a vinegar--oil monomial only when its second index is exactly
\code{oil\_var\_idx}.  The second-oil predicate repeats the same construction
with prefix length 97.  Algebraically, the generated layers are
\begin{align}
 g_j(w,y)&=q_j(w)+y_j\lambda_j(w),
     &w&\in K^{83}, &1\leq j\leq14, \label{eq:o1}\\
 h_k(a,z)&=r_k(a)+z_k\mu_k(a),
     &a&\in K^{97}, &1\leq k\leq40. \label{eq:o2}
\end{align}
All coefficients in $q_j,r_k,\lambda_j,\mu_k$ are sampled before the
forbidden monomials are zeroed.  With overwhelming probability, the displayed
linear forms are nonzero and independent within each layer.

The diagonal layer already has one new variable per equation.  Its $i$th
output is
\begin{equation}
 f_i(u,d_1,\ldots,d_i)
   =a_i(u,d_1,\ldots,d_{i-1})
    +d_i b_i(u,d_1,\ldots,d_{i-1}),
       \qquad 1\leq i\leq16.                         \label{eq:diag}
\end{equation}
Equations \eqref{eq:o1}--\eqref{eq:diag} show that the entire submitted
central map is a 70-step triangular system with one new variable per output,
although the 14- and 40-equation blocks omit permitted dependencies on earlier
variables in their own blocks.

\subsection{Rank-one derivative lines}

Let $e_k$ be the $k$th coordinate vector of the last oil space.  Equation
\eqref{eq:o2} gives
\[
  J^F_{e_k}(a',z')=(0,\ldots,0,\mu_k(a'),0,\ldots,0).
\]
Thus $J^F_{e_k}$ has rank one whenever $\mu_k\neq0$, and its image is the
$k$th second-oil output line.  Also $F(e_k)=0$.

\begin{proposition}[Public oil-line signal]
Let $x_k=T^{-1}e_k$ and let $q_k$ be the image under $S$ of the $k$th
second-oil output line.  Then
\[
       P(x_k)=0,\qquad \rank J^P_{x_k}=1,
       \qquad \im J^P_{x_k}=q_k.
\]
Consequently the projective lines $\langle x_k\rangle$ are detectable from
the public key, their span is the hidden space $O_2$, and the span of their
derivative images is $Q_2$.
\end{proposition}

Inside $O_2$, a vector supported on $s$ coordinate lines has derivative rank
$s$ provided the corresponding $\mu_k$ are independent.  Hence the rank-one
points in $\Proj(O_2)$ are precisely the 40 coordinate lines with overwhelming
probability.  The same argument applies to the 14 first-oil lines after
quotienting by $(O_2,Q_2)$.

\section{Projective multipencil recovery}

Directly enumerating $\Proj(K^{137})$ is impossible.  We instead convert the
rank-one condition into a small search over output ratios.

Choose evaluation points $a_1,\ldots,a_t\in K^n$ and define public linear maps
\[
       L_i:K^n\longrightarrow K^m,\qquad L_i(x)=B_P(x,a_i).
\]
If $\rank J_x=1$, then all $L_i(x)$ lie on one output line.  Therefore there
are $u\in K^m$ and a projective label
$c=(c_1:\cdots:c_t)\in\Proj^{t-1}(K)$ such that $L_i(x)=c_i u$.
Choose a pivot $p$ with $c_p\neq0$.  Then $x$ lies in the kernel of
\begin{equation}
 A_c=\begin{bmatrix}
       c_pL_1-c_1L_p\\[-2pt]
       \vdots\\[-2pt]
       c_pL_t-c_tL_p
     \end{bmatrix}_{i\neq p}.                         \label{eq:Ac}
\end{equation}
For each label $c$, we compute $\ker A_c$ and retain only vectors whose full
public derivative has rank one.  A final check requires $P(x)=0$.

\begin{lemma}[Multipencil completeness]
Let $x\neq0$ satisfy $\rank J_x=1$.  If not all $J_x(a_i)$ vanish, then the
projective label $c=(J_x(a_1):\cdots:J_x(a_t))$ is well defined after
identifying the one-dimensional image of $J_x$ with $K$, and
$x\in\ker A_c$.
\end{lemma}

\begin{proof}
Write $J_x(y)=\ell_x(y)u$ for a nonzero $u$.  Take
$c_i=\ell_x(a_i)$.  Every row block in \eqref{eq:Ac} annihilates $x$ because
$c_pL_i(x)-c_iL_p(x)=c_pc_i u-c_ic_pu=0$.
\end{proof}

For random evaluations, the excluded event has probability $q^{-t}$ for a
fixed rank-one line.  If a run finds fewer than the expected number of
independent lines, it is simply repeated with fresh evaluations.  Every output
is checked by the full derivative, so this is a Las Vegas procedure.

\subsection{Second oil layer}

At the first stage $(n,m)=(137,70)$.  Three evaluations suffice because
$2m=140\geq137$.  There are only
\[
       |\Proj^2(\F_{16})|=16^2+16+1=273
\]
labels.  For each label we reduce a $140\times137$ matrix.  The verified 40
lines span $O_2$; one nonzero column of each rank-one derivative supplies its
image line, and these span $Q_2$.  Public checks give
\[
   \dim O_2=\dim Q_2=40,\quad P(O_2)=0,quad
   B_P(K^{137},O_2)\subseteq Q_2.
\]
The quotient therefore has dimensions
\[
  \overline P_2:K^{97}\longrightarrow K^{30}.
\]

\subsection{First oil layer and the \texorpdfstring{$3+2$}{3+2} split}

For $(n,m)=(97,30)$, five evaluations are needed because
$3m<97\leq4m$.  A direct projective search would have
$(16^5-1)/15=69{,}905$ labels and reduce a $120\times97$ matrix for each one.
We use a $3+2$ split.

First enumerate the 273 projective labels for the first three evaluations.
The first two proportionality relations leave a typical kernel $K_c$ of
dimension
\[
       97-2\cdot30=37.
\]
For every such $K_c$, enumerate the two remaining affine ratios
$(a,b)\in K^2$.  The last two relations are then a $60\times37$ system on
$K_c$.  This requires
\[
       273\cdot16^2=69{,}888
\]
small reductions.  The 17 projective labels whose first three entries vanish
are intentionally omitted.  A desired line is omitted only when three random
evaluations all lie in the kernel of its derivative, an event of probability
$16^{-3}$.  Restarting if fewer than 14 lines are found makes the algorithm
Las Vegas.

The resulting 14 verified lines span $O_1/O_2$ and their derivative images
span $Q_1/Q_2$.  Quotienting again produces
\[
       P_D:K^{83}\longrightarrow K^{16}.
\]

\section{Recovering the diagonal flag}

The quotient $P_D$ is linearly equivalent to the 16 equations in
\eqref{eq:diag}.  Let
\[
       \{0\}=J_0\subset J_1\subset\cdots\subset J_{16}=K^{16}
\]
be the hidden output flag, where $J_\ell$ is the span of the first $\ell$
diagonal forms.  Every form in $J_\ell$ uses at most $v+\ell=67+\ell$
essential variables, so
\begin{equation}
       J_\ell\subseteq
       \mathcal R_{67+\ell}
       \defeq\{c\in K^{16}:\essrank(c\cdot P_D)\leq67+\ell\}.
                                                               \label{eq:rankvar}
\end{equation}

For random coefficients, $J_\ell$ is the unique codimension-one linear
component of
$J_{\ell+1}\cap\mathcal R_{67+\ell}$.  Individual points outside the hidden
hyperplane can be accidentally rank deficient, so accepting a point from one
rank test is unsafe.  We instead grow a linear space while checking projective
line closure.  A candidate $x$ is admitted only when, for random confirmed
$z$ and every $\lambda\in K^*$,
\[
          x+\lambda z\in\mathcal R_{67+\ell}.
\]
After a hyperplane has dimension $\ell$, 32 additional random combinations
are tested.  If a check fails, that descent is restarted.  Beginning with
$J_{16}=K^{16}$ and descending to $J_1$ recovers the complete output flag.

Choose ordered forms $p_1,\ldots,p_{16}$ such that
$\Span(p_1,\ldots,p_\ell)=J_\ell$.  Their common polar radicals reveal the
input-tail flag
\begin{equation}
       K_\ell=\bigcap_{i=1}^{\ell}\rad(B_{p_i}),
       \qquad \dim K_\ell=16-\ell.                    \label{eq:tail}
\end{equation}
The one-dimensional quotients $K_{\ell-1}/K_\ell$ give directions for the
$\ell$th new variable.  Both the dimensions in \eqref{eq:tail} and all rank
bounds in \eqref{eq:rankvar} are public, efficiently checkable certificates
of the recovered equivalent flag.

\section{Public-map inversion and forgery}

Let $t\in K^{70}$ be an arbitrary target.  Project it successively to the
30- and 16-dimensional quotient codomains.  We invert in the opposite order.

\subsection{Inverting the diagonal quotient}

The first ordered form ignores the 15-dimensional tail $K_1$.  Choose a
random point in a fixed 68-dimensional complement until
$p_1(x)=t_1$; this takes about 16 trials.  For
$\ell=2,\ldots,16$, choose a representative
$w_\ell\in K_{\ell-1}\setminus K_\ell$ and enumerate
$a\in\F_{16}$ until
\[
       p_\ell(x+a w_\ell)=t_\ell.
\]
Earlier equations remain unchanged because $w_\ell$ lies in their common
polar radical and in their zero tail.  If a layer has no solution, restart
from a new first-layer point.  The final public evaluation certifies the
quotient preimage.

\subsection{Linear lifts}

Let $x$ be a representative of the diagonal-quotient solution in $K^{97}$.
For $o\in O_1/O_2$, the quotient map satisfies
\[
  \overline P_2(x+o)=\overline P_2(x)+B_{\overline P_2}(x,o),
\]
because the recovered oil space is a zero subspace.  The residual lies in
$Q_1/Q_2$, so its cancellation is a linear system in the 14 coordinates of
$o$.  If the system is singular or inconsistent, use another diagonal
preimage.

Lift the resulting point to $K^{137}$ and repeat the same operation with
$o\in O_2$.  This time the linear system has 40 unknowns and a residual in
$Q_2$.  A successful solution $s$ is checked by evaluating the full public
map and requiring $P(s)=t$.

To forge a signature on a chosen fresh message $M$, choose a 16-byte salt
$\rho$, compute the submission's target
\[
       t=H(H(M)\mathbin\Vert\rho)\in\F_{16}^{70},
\]
and apply the inversion above.  Pack $s\in\F_{16}^{137}$ in 69 bytes and
append $\rho$.  This is an 85-byte signature in the submitted format.

\begin{theorem}[Conditional correctness]
Assume that the two multipencil searches recover the stated rank-one line
sets, the diagonal rank filtration returns flags satisfying
\eqref{eq:rankvar}--\eqref{eq:tail}, and the two lift systems have a solution
for one sampled quotient preimage.  Then the algorithm outputs a signature
accepted by the submitted VDOO-128 verifier for the chosen message and salt.
\end{theorem}

\begin{proof}
The quotient lemma shows that each recovered oil pair defines the induced
quadratic map used by the next stage.  The tail directions preserve all
previous diagonal equations, and exhaustive enumeration in $\F_{16}$ sets the
next equation; hence the diagonal output is correct.  Each lift uses the
identity $P(x+o)=P(x)+B_P(x,o)$ on a zero oil space and solves exactly for the
remaining residual.  Therefore the final vector satisfies
$P(s)=H(H(M)\Vert\rho)$, which is precisely the submitted verification
condition.
\end{proof}

All hypotheses are checked from the public key.  A failed random choice causes
a restart rather than an incorrect output.  Final verifier acceptance is an
end-to-end certificate independent of the structural interpretation.

\section{Complexity}

Table~\ref{tab:work} gives the fixed search sizes for VDOO-128.  The costs are
polynomial except for enumerations over very small projective spaces whose
sizes are shown explicitly.  There is no term near $2^{128}$.

\begin{table}[ht]
\centering
\caption{Main public-key work for submitted VDOO-128.}
\label{tab:work}
\small
\begin{tabular}{@{}p{38mm}p{43mm}p{55mm}@{}}
\toprule
Stage & Search & Dominant linear algebra \\
\midrule
Recover $O_2,Q_2$ & 273 labels in $\Proj^2(\F_{16})$
  & kernels of $140\times137$ matrices \\
Recover $O_1/O_2,Q_1/Q_2$ & $273\cdot256=69{,}888$ split labels
  & 273 kernels of $60\times97$, then kernels of $60\times37$ \\
Recover diagonal flag & about $5.3\times10^3$ sampled projective rank tests
  & essential rank of forms on at most 83 variables \\
Invert one target & about $16$ trials for the first form, plus retries
  & 14- and 40-unknown linear lifts \\
\bottomrule
\end{tabular}
\end{table}

The projective searches can be parallelized label by label.  Our proof-of-
concept does not do so.  It uses byte lookup tables for $\F_{16}$ arithmetic
and vectorized row operations; it does not use compiled finite-field kernels,
SIMD, GPUs, or multiple cores.  The reported timings should therefore be read
as a practical upper bound rather than an optimized record.

\section{Experiments}

\subsection{Environment and methodology}

We built the submitted VDOO-128 reference sources from NGCC harness commit
\code{32f381e}.  Experiments ran on an
Intel Core i5-6500 at 3.20~GHz (four physical cores), with 15~GiB RAM, Python
3.13.13, and NumPy 2.5.3.  The attack process is single-threaded.

For the independent-key trial, we initialized the implementation's private
DRNG with the 48-byte seed $00,01,\ldots,2f$ before key generation.  The secret
key returned by key generation was discarded.  The attack process received
only the serialized public key.  This bypasses the already reported
zero-initialized-RNG behavior and tests a different, independently sampled
central map and hiding transformations.

\subsection{Results}

\begin{table}[ht]
\centering
\caption{Uncached end-to-end attack on the independently seeded key.  ``Online''
includes quotient inversion, both lifts, signature packing, and verification.}
\label{tab:timings}
\footnotesize
\begin{tabular}{@{}p{29mm}rrrrrl@{}}
\toprule
Trial & $O_2$ & $O_1$ & flag/checks & online & total & result \\
\midrule
Independent seeded key & 25.0 s & 380.4 s & 201.8 s & 78.0 s & 693.2 s & accept \\
\bottomrule
\end{tabular}
\end{table}

For both keys the first multipencil returned exactly 40 projective rank-one
lines spanning a 40-dimensional public zero subspace, and their images spanned
40 dimensions.  The quotient multipencil returned exactly 14 lines spanning
a 14-dimensional zero subspace and 14 output dimensions.  The diagonal stage
recovered all 15 proper output-flag subspaces and the tail dimensions
$15,14,\ldots,0$.  Every stage was checked using only the public map.

The first complete structural run took 600.9 seconds to invert an arbitrary
target.  Reusing the public-key structure, a subsequent fresh-message forgery
was accepted in 25.5 seconds of online inversion.  The independent seeded-key
run started without cached structure and ended with acceptance by the same
submitted verification function.  Peak resident memory remained below the
available 15~GiB; the persistent public-key attack state occupies less than
one megabyte.

\subsection{Accepted independent-key forgery}

The independently seeded experiment forged the following message without a
signing query:
\begin{quote}
\small\ttfamily Independent-key public-key-only\\
structural forgery against submitted VDOO-128
\end{quote}
The submitted verifier returned zero (accept).  Flipping one bit of the
message or one bit of the signature made it return $-1$ (reject).  The binary
artifacts have the following SHA-256 digests:
\begin{center}
\small
\begin{tabular}{@{}ll@{}}
\toprule
Object & SHA-256 \\
\midrule
Public key & \code{f0bf3e7faaa200c14c8b37c12d4a0d7195d2ff6c0ce8081e29ea303a4b3d91c6} \\
Message & \code{31b942f4df86978090308e2bc13a54379dd0333c66d5210831d4ddf49ceeb575} \\
Signature & \code{575c13a9f988e6210c562c2cab861dc13228878c1aa45d7d7357f6188ea5b08c} \\
\bottomrule
\end{tabular}
\end{center}
The 85-byte signature, including its final 16-byte salt, is
\begin{quote}
\ttfamily\small\raggedright
9f3606a0748596e7bf00a719fd1234dbb562a86fa6f465e377f81d1f6b8f8d63\\
f28b5577b42cc809d5d76d9fdbb77c7b6533b561a2aef559b986dbc4972c1311\\
3ba7d2e00f33d253ff531ba22f82951e499884ed23
\end{quote}
These values are included to make the acceptance claim independently
checkable without relying on any attack implementation.

\section{Why the attack survives an RNG repair}

The public RNG failure allows an even simpler secret-key reproduction attack
\cite{ngcc-report}.  It is not used here.  The signal exploited in this paper
is the public tensor identity
\[
       \rank\bigl(y\mapsto B_P(x_k,y)\bigr)=1
\]
for every hidden matching-oil direction.  This identity follows from which
coefficients are set to zero after random sampling.  Reseeding changes the
nonzero coefficients and hiding transformations, but it preserves the
identity.  The independently seeded forgery confirms this distinction.

Likewise, the attack needs no fault, timing trace, malformed key, reused
signature randomness, chosen-target oracle, or known signature.  Its input is
a well-formed public key.  Since it forges a fresh chosen message in the
key-only setting, it directly violates existential unforgeability and is
stronger than an attack requiring chosen-message access.

\section{Remediation and design lesson}

The immediate repair is to generate the full vinegar--oil block specified for
each oil equation.  For the first layer, output $r$ must permit
\[
  0\leq i<83,\qquad 83\leq j<97,
\]
for every $j$, rather than only $j=83+r$.  For the second layer, it must permit
\[
  0\leq i<97,\qquad 97\leq j<137
\]
for every $j$.  Test vectors and key-size calculations must then be regenerated
and the repaired distribution reanalyzed.  The RNG and side-channel findings
require separate fixes.

Unit tests should test structural distribution properties, not only
sign--verify correctness.  In particular, for random secret-coordinate oil
directions in a dense UOV layer, the derivative image should span many output
coordinates.  A test that counts the public rank-one derivative locus after a
small hidden toy instantiation would have exposed the one-variable mask while
ordinary known-answer tests continued to pass.

The broader lesson is that a coefficient filter is part of the cryptographic
design.  Replacing a dense UOV linear system by per-equation matching columns
does more than optimize storage or signing: it changes the isomorphism class
of the public quadratic tensor and exposes invariant rank-one points.  Hiding
linear transformations cannot erase those tensor ranks.

\section{Conclusion}

The submitted VDOO-128 implementation can be forged from its public key in
minutes on an ordinary laptop-class machine.  The attack recovers rank-one
oil directions, forms two public quotients, reconstructs an equivalent
diagonal flag, and inverts a fresh hash target.  A signature produced for an
independently seeded key is accepted by the submitted verifier.  The attack is
caused by a sparse coefficient mask that selects only the oil variable matching
each output equation.  Restoring the dense UOV layers in the normative
specification removes this particular rank-one mechanism, but requires new
test vectors and fresh cryptanalysis.

\appendix
\section{Attack pseudocode}

\paragraph{Rank-one multipencil.}
Given the polar map $B:K^n\times K^n\to K^m$ and an evaluation count $t$:
\begin{enumerate}[leftmargin=*,itemsep=1pt]
  \item sample $a_1,\ldots,a_t\leftarrow K^n$ and set
        $L_i(x)=B(x,a_i)$;
  \item enumerate $c\in\Proj^{t-1}(K)$, choose a nonzero pivot $c_p$, and
        form $A_c$ as in \eqref{eq:Ac};
  \item enumerate the projective points of each small $\ker A_c$;
  \item retain $x$ precisely when $\rank J_x=1$ and $P(x)=0$;
  \item normalize projective representatives and remove duplicates;
  \item restart with new $a_i$ if the verified span has the wrong dimension.
\end{enumerate}

For the $3+2$ variant, enumerate a projective label for $L_1,L_2,L_3$,
compute the kernel of the first two relations, restrict $L_p,L_4,L_5$ to that
kernel, and enumerate the final two affine ratios.

\paragraph{Full inversion.}
\begin{enumerate}[leftmargin=*,itemsep=1pt]
  \item use the three-evaluation multipencil to recover $O_2,Q_2$;
  \item construct $P_2:K^{137}/O_2\to K^{70}/Q_2$;
  \item use the five-evaluation split multipencil to recover
        $O_1/O_2,Q_1/Q_2$;
  \item construct $P_D:K^{83}\to K^{16}$;
  \item recover $J_{15},\ldots,J_1$ by descending essential-rank
        hyperplanes and derive the tail radicals $K_1,\ldots,K_{16}$;
  \item invert the projected target through the 16 one-dimensional layers;
  \item solve the 14-variable first-oil lift and the 40-variable second-oil
        lift, restarting the diagonal inversion after an inconsistent lift;
  \item require full public evaluation $P(s)=t$ before returning $s$.
\end{enumerate}

\small\sloppy
\begin{thebibliography}{99}

\bibitem{vdoo-original}
A.~Ganguly, A.~Karmakar, and N.~Saxena.
\newblock VDOO: A short, fast, post-quantum multivariate digital signature
  scheme.
\newblock \emph{Cryptology ePrint Archive}, Report 2023/1925, 2023.
\newblock \url{https://eprint.iacr.org/2023/1925}.

\bibitem{vdoo-spec}
S.~Adhikary, C.~Bravo, P.~Coscojuela-Escanilla, A.~Ganguly, A.~Karmakar,
S.~Kundu, D.~Pal, L.~Perret, and N.~Saxena.
\newblock VDOO: Vinegar Diagonal Oil Oil, Specification v1.1.
\newblock NGCC submission specification, 30 June 2026.
\newblock \url{https://cdn.jsdelivr.net/gh/ngcc-dev/ngcc-harness@main/sign-33/sign-33-spec.pdf}.

\bibitem{ngcc-source}
NGCC reproduction harness.
\newblock VDOO (sign-33) reference sources at commit \code{32f381e},
  24 September 2026.
\newblock \href{https://github.com/ngcc-dev/ngcc-harness/tree/32f381e40ae2eb35c83a454e8e59b17f44c26c7d/sign-33}{Archived source snapshot}.

\bibitem{ngcc-report}
NGCC reproduction project.
\newblock VDOO (sign-33) security report, accessed 25 September 2026.
\newblock \url{https://ngcc.dev/reports/sign-33.html}.

\bibitem{wolf-sts}
C.~Wolf, A.~Braeken, and B.~Preneel.
\newblock Efficient cryptanalysis of RSE(2)PKC and RSSE(2)PKC.
\newblock \emph{Cryptology ePrint Archive}, Report 2004/237, 2004.
\newblock \url{https://eprint.iacr.org/2004/237}.

\bibitem{beullens-uov}
W.~Beullens.
\newblock Improved cryptanalysis of UOV and Rainbow.
\newblock \emph{Cryptology ePrint Archive}, Report 2020/1343, 2020;
  EUROCRYPT 2021.
\newblock \url{https://eprint.iacr.org/2020/1343}.

\bibitem{beullens-rainbow}
W.~Beullens.
\newblock Breaking Rainbow takes a weekend on a laptop.
\newblock \emph{Cryptology ePrint Archive}, Report 2022/214, 2022.
\newblock \url{https://eprint.iacr.org/2022/214}.

\end{thebibliography}

\end{document}
